\documentclass[10pt,a4paper]{amsart}
\usepackage{amsmath,amssymb,amsthm,graphicx}
\usepackage{verbatim}
\newenvironment{reference}{\comment}{\endcomment}
%\newenvironment{reference}{}{}
\newenvironment{slogan}{\comment}{\endcomment}
\newenvironment{history}{\comment}{\endcomment}

% For commutative diagrams we use Xy-pic
\usepackage[all]{xy}

% We use 2cell for 2-commutative diagrams.
\xyoption{2cell}
% Theorem environments.
%
\theoremstyle{plain}
\newtheorem{theorem}[subsection]{Theorem}
\newtheorem{proposition}[subsection]{Proposition}
\newtheorem{lemma}[subsection]{Lemma}

\theoremstyle{definition}
\newtheorem{definition}[subsection]{Definition}
\newtheorem{example}[subsection]{Example}
\newtheorem{exercise}[subsection]{Exercise}
\newtheorem{situation}[subsection]{Situation}

\theoremstyle{remark}
\newtheorem{remark}[subsection]{Remark}
\newtheorem{remarks}[subsection]{Remarks}

\numberwithin{equation}{subsection}

% Macros
%
\def\lim{\mathop{\mathrm{lim}}\nolimits}
\def\colim{\mathop{\mathrm{colim}}\nolimits}
\def\Spec{\mathop{\mathrm{Spec}}}
\def\Hom{\mathop{\mathrm{Hom}}\nolimits}
\def\Ext{\mathop{\mathrm{Ext}}\nolimits}
\def\SheafHom{\mathop{\mathcal{H}\!\mathit{om}}\nolimits}
\def\SheafExt{\mathop{\mathcal{E}\!\mathit{xt}}\nolimits}
\def\Sch{\mathit{Sch}}
\def\Mor{\mathop{\mathrm{Mor}}\nolimits}
\def\Ob{\mathop{\mathrm{Ob}}\nolimits}
\def\Sh{\mathop{\mathit{Sh}}\nolimits}
\def\NL{\mathop{N\!L}\nolimits}
\def\CH{\mathop{\mathrm{CH}}\nolimits}
\def\proetale{{pro\text{-}\acute{e}tale}}
\def\etale{{\acute{e}tale}}
\def\QCoh{\mathit{QCoh}}
\def\Ker{\mathop{\mathrm{Ker}}}
\def\Im{\mathop{\mathrm{Im}}}
\def\Coker{\mathop{\mathrm{Coker}}}
\def\Coim{\mathop{\mathrm{Coim}}}

% Boxtimes
%
\DeclareMathSymbol{\boxtimes}{\mathbin}{AMSa}{"02}

%
% Macros for moduli stacks/spaces
%
\def\QCohstack{\mathcal{QC}\!\mathit{oh}}
\def\Cohstack{\mathcal{C}\!\mathit{oh}}
\def\Spacesstack{\mathcal{S}\!\mathit{paces}}
\def\Quotfunctor{\mathrm{Quot}}
\def\Hilbfunctor{\mathrm{Hilb}}
\def\Curvesstack{\mathcal{C}\!\mathit{urves}}
\def\Polarizedstack{\mathcal{P}\!\mathit{olarized}}
\def\Complexesstack{\mathcal{C}\!\mathit{omplexes}}
% \Pic is the operator that assigns to X its picard group, usage \Pic(X)
% \Picardstack_{X/B} denotes the Picard stack of X over B
% \Picardfunctor_{X/B} denotes the Picard functor of X over B
\def\Pic{\mathop{\mathrm{Pic}}\nolimits}
\def\Picardstack{\mathcal{P}\!\mathit{ic}}
\def\Picardfunctor{\mathrm{Pic}}
\def\Deformationcategory{\mathcal{D}\!\mathit{ef}}
\usepackage{cleveref}
\crefname{originalreference}{original reference}{original references}
\Crefname{originalreference}{Original reference}{Original references}
\makeatletter
\newlabel{lemma-localization}{{096V}{source}{}{sourcepointer-1}{}}
\newlabel{lemma-localization@cref}{{[originalreference][0][]096V}{[1][0][]source}{}{}{}}
\makeatother
\crefname{originalsection}{original section}{original sections}
\Crefname{originalsection}{Original section}{Original sections}
\makeatletter
\newlabel{sites-modules-section-localize}{{03DH}{source}{}{sourcepointer-2}{}}
\newlabel{sites-modules-section-localize@cref}{{[originalsection][0][]03DH}{[1][0][]source}{}{}{}}
\makeatother
\title{Vanishing of the extension by zero along a weakly etale morphism over an affine object all of whose points specialize into a disjoint closed part}
\author{The Stacks Project Authors}
\date{stacks-project:proetale.tex:local-snapshot-20261012; GNU FDL 1.2 or later}
\begin{document}
\maketitle
\noindent\textit{Editorial note.} All delivered mathematical text is the author's own, copied from the pinned Stacks Project chapter file proetale.tex; no statement or proof step was rewritten, reordered or removed.\par\medskip
\noindent\textit{Editorial note.} The counted proof is the author's complete argument: the criterion that the sections of the extension by zero are computed as a colimit over coverings whose members admit no morphism into U, the construction of the refining covering by localizing at the part over Z, and the closing specialization argument excluding morphisms into U.\par\medskip
\noindent\textit{Editorial note.} The label lemma-jshriek-zero names the counted statement, and the proof environment immediately following it in the source is delivered verbatim as its proof.\par\medskip
\noindent\textit{Editorial note.} The counted statement is carried by the author's own lemma environment in the chapter source.\par\medskip
\noindent\textit{Editorial note.} Pre-open gate: python operations/agent-authoring/tools/precard\_gate.py --source proetale.tex --label lemma-jshriek-zero --start 5306 --end 5331 --local-id 3599 returned RESULT=OK (L3599 free, occupancy FREE, public interval clear). Candidate taken from sources/stacks-proetale/candidates-halfA.json.\par\medskip
\noindent\textit{Editorial note.} This card closes the first numbering segment L3400-L3599; per LANE\_BRIEF section 4 the following cards use L28000-L29999.\par\medskip
\noindent\textit{Editorial note.} The statement and proof are copied from the pinned chapter file; the cited inputs are kept as external pointers with their canonical Stacks tags.\par\medskip
\noindent\textit{Allowed standard prerequisites (editorial).}\begin{itemize}
\item Construction of the extension by zero by refining coverings. Modules on Sites, the section constructing the extension by zero together with the description of its sections as a colimit over coverings whose members admit no morphism to the open part. Supplies the criterion that the sections of the extension by zero vanish once the covering can be refined to members having no morphism into U over X.
\item The localization construction attached to a locally closed subscheme. Same chapter, the lemma producing a ring A\_Z\textasciicircum{}tilde with spectrum the points specializing to a locally closed subscheme Z of the form D(f) cap V(I). Supplies the affine schemes W\_i = Spec(A\_\{i,Z\_i\}\textasciicircum{}tilde) whose images contain all points of V\_Z and hence all points of V.
\end{itemize}\medskip
\ifdefined\hypertarget\hypertarget{sourcepointer-1}{}\fi
\noindent\textit{Source reference.} Original reference 096V: sources/stacks-proetale/pinned-original-source.tex, source lines 710--710. Prerequisite pointer invoked at proof line(s) 5324; it names a standard fact of the same chapter that the counted argument uses, and every step of the counted proof built on it is retained verbatim in the delivered body.\par\medskip
\ifdefined\hypertarget\hypertarget{sourcepointer-2}{}\fi
\noindent\textit{Source reference.} Original section 03DH: sources/stacks-chow/tags.txt, source lines 4420--4420. Prerequisite pointer invoked at proof line(s) 5319; it names a standard fact of a sibling chapter that the counted argument uses, and every step of the counted proof built on it is retained verbatim in the delivered body.\par\medskip
\begin{lemma}
\label{lemma-jshriek-zero}
Let $j : U \to X$ be a weakly \'etale morphism of schemes.
Let $i : Z \to X$ be a closed immersion such that $U \times_X Z = \emptyset$.
Let $V \to X$ be an affine object of $X_\proetale$ such that every point
of $V$ specializes to a point of $V_Z = Z \times_X V$.
Then $j_!\mathcal{F}(V) = 0$ for all abelian sheaves on $U_\proetale$.
\end{lemma}

\begin{proof}
Let $\{V_i \to V\}$ be a pro-\'etale covering. The lemma follows if we
can refine this covering to a covering where the members have no
morphisms into $U$ over $X$ (see construction of $j_!$ in
Modules on Sites, Section \ref{sites-modules-section-localize}).
First refine the covering to get a finite covering with $V_i$ affine.
For each $i$ let $V_i = \Spec(A_i)$ and let $Z_i \subset V_i$ be the
inverse image of $Z$.
Set $W_i = \Spec(A_{i, Z_i}^\sim)$ with notation as in
Lemma \ref{lemma-localization}.
Then $\coprod W_i \to V$ is weakly \'etale and the image contains all
points of $V_Z$. Hence the image contains all points of $V$ by
our assumption on specializations. Thus $\{W_i \to V\}$ is a
pro-\'etale covering refining the given one. But each point in $W_i$
specializes to a point lying over $Z$, hence there are no morphisms
$W_i \to U$ over $X$.
\end{proof}

\end{document}
